\begin{abstract}
A robot policy can learn from its own attempts using simulator success feedback, without additional demonstrations. We investigate how to sustain this learning as the policy and its training program evolve. In our recursive self-improvement framework, each retained policy generates experience for the next round, and round records guide revisions to the training program. Preregistered experiments with a 7B OpenVLA-OFT policy on two LIBERO suites show that updating the data-generating policy supports continued learning. With a fixed GRPO program, twelve rounds improve success on 1{,}000 fresh scenes by $+7.4$ points $[+5.8, +9.1]$ on LIBERO-Spatial and $+3.5$ $[+1.7, +5.4]$ on LIBERO-Goal. When the loop starts from a behaviour-cloning program that degrades during self-training, a validation gate preserves initial performance and program adaptation restores learning. Continuing under the adapted program then yields a $+7.6$-point advantage $[+5.6, +9.6]$ over tuned one-shot RL on fresh Goal scenes, using 190 versus 230 GPU-hours on average over two seeds, including RL's recorded tuning cost. Code, preregistrations and run records are released.
\end{abstract}
